% ============================================================================
% คู่มือการใช้งานระบบจัดตารางสอบ Exam Scheduler — สไตล์และคำสั่งเฉพาะ
% ใช้ร่วมกับ LuaLaTeX + babel (ไทย) + Sarabun/CommitMono จาก docs/fonts/
% ============================================================================

% --- เรขาคณิตของหน้ากระดาษ (A4 แนวตั้ง) --------------------------------------
\usepackage[left=3cm,right=2.5cm,top=2.5cm,bottom=2.5cm]{geometry}

% --- ฟอนต์ไทยและฟอนต์รหัส/ชื่อไฟล์ ------------------------------------------
% ใช้ Sarabun เป็นฟอนต์เนื้อหา และ CommitMono สำหรับรหัส/ชื่อไฟล์/ตัวเลขเก้าหลัก
\babelfont{rm}[Path=docs/fonts/,Extension=.ttf,UprightFont=Sarabun-Regular,BoldFont=Sarabun-Bold,ItalicFont=Sarabun-Italic,BoldItalicFont=Sarabun-BoldItalic,Scale=1.06,Renderer=HarfBuzz]{Sarabun}
\babelfont{tt}[Path=docs/fonts/,Extension=.otf,UprightFont=CommitMono-400-Regular,BoldFont=CommitMono-700-Regular,ItalicFont=CommitMono-400-Italic,BoldItalicFont=CommitMono-700-Italic,Scale=MatchLowercase,Ligatures={TeXOff,NoContextual,NoDiscretionary},Renderer=HarfBuzz]{CommitMono}
\linespread{1.30}

% --- แพ็กเกจที่ใช้ -----------------------------------------------------------
\usepackage{booktabs,longtable,array,graphicx,tikz,xcolor}
\usepackage{float}
\usepackage{needspace}
% สัญลักษณ์ในตารางเปรียบเทียบ (Sarabun ไม่มีอักขระ U+2713
% จึงใช้เครื่องหมายจากฟอนต์ Zapf Dingbats ของ pifont แทน)
\usepackage{pifont}
\newcommand{\yesmark}{\ding{51}}
\newcommand{\nomark}{\textendash}
\usepackage{ragged2e}
\usepackage{framed}
\usepackage[titles]{tocloft}
\usepackage{url}
\usepackage[unicode,colorlinks=true,linkcolor=black,urlcolor=blue,citecolor=black]{hyperref}

% Flexible spacing at safe word boundaries also works in headings/bookmarks.
\DeclareRobustCommand{\thaiwordbreak}{\hspace{0pt plus 0.1em}}
\pdfstringdefDisableCommands{\def\thaiwordbreak{}}

% Thai word boundaries are supplied by prepare-thai-typesetting.py during the
% build. Disable Babel's character-level Southeast-Asian fallback so a word
% never splits across two lines.
\directlua{Babel.sea_enabled = false}
\hyphenpenalty=10000
\exhyphenpenalty=10000

% --- สารบัญ: จุดนำสายตา (dot leaders) สม่ำเสมอทุกระดับ --------------------------
% cftdotsep มาตรฐาน (4.5) ทำให้จุดห่างเกินในสารบัญรูปซึ่ง caption ยาว
% ลดเหลือ 2.5 ให้จุดถี่สม่ำเสมอทั้งสารบัญ สารบัญรูป และสารบัญตาราง
\renewcommand{\cftdotsep}{2.5}
\renewcommand{\cftsecleader}{\cftdotfill{\cftdotsep}}
\renewcommand{\cftsubsecleader}{\cftdotfill{\cftdotsep}}
\renewcommand{\cftsubsubsecleader}{\cftdotfill{\cftdotsep}}
\renewcommand{\cftfigleader}{\cftdotfill{\cftdotsep}}
\renewcommand{\cfttableader}{\cftdotfill{\cftdotsep}}
\setcounter{tocdepth}{1}
\hypersetup{bookmarksdepth=3}

% --- การแบ่งบรรทัดของรหัส (โค้ด/ชื่อไฟล์) ------------------------------------
% \code/\path กินช่องว่าง ดังนั้นคำสั่งแบบเว้นวรรคต้องใช้ \texttt
% ไม่อนุญาตให้ตัดรหัสกลางคำโดยไม่ตั้งใจ ยกเว้นจุดที่กำหนดไว้
\newcommand{\codebreaks}{\do\.\do\_\do\:\do\/\do\-\do\;}
\newcommand{\code}[1]{\bgroup\def\UrlBreaks{\codebreaks}\path{#1}\egroup}

% --- รูปแบบย่อหน้า -----------------------------------------------------------
\setlength{\parindent}{1.1cm}
\setlength{\parskip}{5pt}
\setlength{\emergencystretch}{3em}
\setlength{\tabcolsep}{4pt}
\renewcommand{\arraystretch}{1.45}
\newcolumntype{L}[1]{>{\raggedright\arraybackslash}p{#1}}
\newcolumntype{R}[1]{>{\raggedleft\arraybackslash}p{#1}}

% Justify prose using stretch distributed over dictionary-approved Thai word
% boundaries. Table columns keep their explicit alignment.
\justifying
\setlength{\parindent}{1.1cm}

% --- คำสั่ง \texttt สำหรับป้ายกำกับภาษาไทย ----------------------------------
% CommitMono มีเฉพาะอักษรละติน จึงไม่สามารถแสดงภาษาไทยได้
% คำสั่ง \texttt ในคู่มือใช้แสดงป้ายกำกับ/ชื่อปุ่มภาษาไทยเป็นหลัก
% จึงนิยามให้ใช้ฟอนต์เนื้อหา (Sarabun) แบบหนาแทน ซึ่งรองรับทั้งไทยและละติน
% ส่วนรหัส/ชื่อไฟล์ที่เป็นอักษรละตินล้วนให้ใช้ \code (CommitMono) แทน
\renewcommand{\texttt}[1]{\textbf{#1}}

% --- การนับหัวข้อและรูป/ตาราง -------------------------------------------------
\counterwithin{figure}{section}
\counterwithin{table}{section}

% ไฟล์แต่ละบทเขียนด้วย \mychapter + \section/\subsection เพื่อให้อ่าน source ง่าย
% แต่ในเอกสารจริงต้องเป็นลำดับ บท > หัวข้อ > หัวข้อย่อย ไม่ใช่รายการต่อเนื่อง 1--82
\let\manualchaptercommand\section
\let\manualsectioncommand\subsection
\let\manualsubsectioncommand\subsubsection
\RenewDocumentCommand{\section}{s m}{%
  \IfBooleanTF{#1}{\manualsectioncommand*{#2}}{\manualsectioncommand{#2}}%
}
\RenewDocumentCommand{\subsection}{s m}{%
  \IfBooleanTF{#1}{\manualsubsectioncommand*{#2}}{\manualsubsectioncommand{#2}}%
}
\newcommand{\mychapter}[1]{%
  \clearpage
  \refstepcounter{section}%
  \phantomsection
  \addcontentsline{toc}{section}{\protect\numberline{\thesection}#1}%
  \markboth{#1}{#1}%
  \noindent{\Large\bfseries \thesection\par}%
  \vspace{0.3em}%
  \noindent{\Huge\bfseries #1\par}%
  \vspace{1.1em}%
}

% เวอร์ชันควบคุมจุดขึ้นบรรทัดใหม่ของชื่อบทเอง โดยให้สารบัญใช้ชื่อเต็มเสมอ
% \mychapterlong{ชื่อในสารบัญ}{ชื่อที่แสดงบนหัวบท (ใส่ \\ ได้)}
\newcommand{\mychapterlong}[2]{%
  \clearpage
  \refstepcounter{section}%
  \phantomsection
  \addcontentsline{toc}{section}{\protect\numberline{\thesection}#1}%
  \markboth{#1}{#1}%
  \noindent{\Large\bfseries \thesection\par}%
  \vspace{0.3em}%
  \noindent{\Huge\bfseries #2\par}%
  \vspace{1.1em}%
}

% --- สไตล์เฉพาะของคู่มือ ------------------------------------------------------
% ขั้นตอนแบบเรียงลำดับ
\newenvironment{steps}{%
  \begin{enumerate}\setlength{\itemsep}{3pt}\setlength{\parskip}{0pt}%
}{%
  \end{enumerate}%
}

% --- บล็อกผลที่ควรเห็น/ข้อควรระวัง (ใช้เป็น environment) -----------------------
% กรอบแถบสีซ้ายด้วย framed leftbar: ไม่มีเยื้องย่อหน้า แถบสูงเท่าเนื้อหา
% ไม่มีช่องว่างบนล่างเกิน และตัดข้ามหน้าได้อย่างถูกต้อง
% expected = แถบเขียว, caution = แถบแดง
\newenvironment{expected}{%
  \par\smallskip\noindent%
  \def\FrameCommand{{\color{green!55!black}\vrule width 3pt}\hspace{8pt}}%
  \setlength{\FrameSep}{0pt}%
  \setlength{\FrameRule}{0pt}%
  \setlength{\topsep}{2pt}\setlength{\partopsep}{0pt}%
  \begin{leftbar}\justifying\setlength{\parindent}{0pt}\noindent%
  {\small\textbf{ผลที่ควรเห็น:} }%
}{%
  \end{leftbar}\par%
}

\newenvironment{caution}{%
  \par\smallskip\noindent%
  \def\FrameCommand{{\color{red!70!black}\vrule width 3pt}\hspace{8pt}}%
  \setlength{\FrameSep}{0pt}%
  \setlength{\FrameRule}{0pt}%
  \setlength{\topsep}{2pt}\setlength{\partopsep}{0pt}%
  \begin{leftbar}\justifying\small\setlength{\parindent}{0pt}\noindent%
  \textcolor{red!70!black}{\textbf{ข้อควรระวัง:} }%
}{%
  \end{leftbar}\par%
}

% หมายเหตุ = แถบเทากลาง (ข้อมูลเพิ่มเติมที่ไม่ใช่ความเสี่ยง)
% แยกจาก caution (แดง) เพื่อให้คำเตือนที่สำคัญจริงยังโดดเด่น
\newenvironment{note}{%
  \par\smallskip\noindent%
  \def\FrameCommand{{\color{black!45}\vrule width 3pt}\hspace{8pt}}%
  \setlength{\FrameSep}{0pt}%
  \setlength{\FrameRule}{0pt}%
  \setlength{\topsep}{2pt}\setlength{\partopsep}{0pt}%
  \begin{leftbar}\justifying\small\setlength{\parindent}{0pt}\noindent%
  \textcolor{black!60}{\textbf{หมายเหตุ:} }%
}{%
  \end{leftbar}\par%
}

% คำอธิบายภาพ — ใช้ \figcaption ใน figure แทน
\newcommand{\figcaption}[1]{\caption{#1}}

% คำศัพท์/คำย่อ
\newcommand{\gloss}[2]{\textbf{#1}: #2}

% การอ้างอิงข้ามบท: \seealso{หมายเลขขั้นตอน/หัวข้อ}
\newcommand{\seealso}[1]{\par\smallskip{\small ดูเพิ่มเติม: #1}\par}

% ชื่อไฟล์ที่ใช้ในคู่มือ (ไม่ให้ตัดกลาง)
% ใช้ \texttt (Sarabun หนา) แทน \code เพื่อรองรับชื่อไฟล์ภาษาไทย
\newcommand{\file}[1]{\texttt{#1}}

% รหัสวิชาเก้าหลัก — ใช้ \texttt (Sarabun หนา) ซึ่งปลอดภัยในคำอธิบายภาพ
% เนื่องจากรหัสวิชาเป็นตัวเลขล้วน ไม่ต้องใช้ฟอนต์โมโน
\newcommand{\course}[1]{\texttt{#1}}

% คำสั่งสำหรับแทรกรูปพร้อมคำอธิบาย
% \manualfigure{ไฟล์}{ความกว้าง}{คำอธิบาย}
% ใช้ตำแหน่ง [H] เพื่อให้ภาพหน้าจอประกอบขั้นตอนอยู่ติดกับข้อความที่สอนเสมอ
% ไม่ลอยไปรวมกันท้ายหน้า/ท้ายบทจนเกิดหน้าที่มีแต่รูปกับพื้นที่ว่าง
\newcommand{\manualfigure}[3]{%
  \begin{figure}[H]\centering
    \includegraphics[width=#2]{#1}
    \caption{#3}
  \end{figure}%
}

% --- เมตาดาตาของ PDF --------------------------------------------------------
\hypersetup{%
  pdftitle={คู่มือการใช้งานระบบจัดตารางสอบ Exam Scheduler},
  pdfauthor={Exam Scheduler project},
  pdflang={th},
  pdfsubject={คู่มือการใช้งานสำหรับบุคลากรจัดตารางสอบ},
  pdfkeywords={จัดตารางสอบ, Exam Scheduler, คู่มือ, ไทย}
}
